\documentclass[a4paper,12pt,twoside,openright,titlepage]{article}

%Additional packages
\input{styles/packages}

% Listing style
\input{styles/lstset}

% Uncomment for production
% \syntaxonly

% Style
\pagestyle{headings}

% Define document
\author{D. Leeuw}
\title{PowerShell: Opdracht werken met variabelen}
\date{\today\\
0.0.0 \\
\vfill
\raggedright
\copyright\ 2026 Dennis Leeuw\\
\input{styles/licentie-titlepage}}


\begin{document}
\selectlanguage{dutch}

\maketitle

%%%%%%%%%%%%%%%%%%%
%%% Introductie %%%
%%%%%%%%%%%%%%%%%%%

\section{Over dit Document}
\subsection{Leerdoelen}
\subsection{Voorkennis}

%%%%%%%%%%%%%%%%%
%%% De inhoud %%%
%%%%%%%%%%%%%%%%%

\section{De situatie}
\input{src/powershell-opdracht-datatypes_situatie}

\section{Kennismaken met variabelen}
\input{src/powershell-opdracht-datatypes_variabelen}

\section{Gebruikers invoer}
\input{src/powershell-opdracht-datatypes_userinput}

\section{Rekenen met variabelen}
\input{src/powershell-opdracht-datatypes_rekenen}

\section{Maak een Systemoverzicht}
\input{src/powershell-opdracht-datatypes_overzicht}

\section{Extra informatie}
\input{src/powershell-opdracht-datatypes_extra}

%%%%%%%%%%%%%%%%%%%%%
%%% Index and End %%%
%%%%%%%%%%%%%%%%%%%%%
%\backmatter
\printindex
\end{document}

%%% Last line %%%
